\documentclass[12pt]{article}
\usepackage[T1]{fontenc}
\usepackage[utf8]{inputenc}
\usepackage{lmodern}
\usepackage{textcomp}
\usepackage{enumitem}
\usepackage{geometry}
\usepackage{graphicx}
\usepackage{amsmath, amssymb}
\geometry{margin=1in}
\begin{document}
\begin{center}
Economics - Undergraduate Level Assessment 15 \
Total Marks: 40 \
Answer all questions.
\end{center}

\section*{Multiple Choice (10 marks)}
\begin{enumerate}[label=\arabic*.]
\item In the last 20 years firms that produce cameras have begun to produce fewer 35-mm cameras and more digital cameras. This trend is an example of
\begin{enumerate}[label=(\alph*)]
    \item the market system answering the question of "what" cameras should be produced.
    \item the market system answering the question of "when" cameras should be produced.
    \item the market system answering the question of "why" cameras should be produced.
    \item the influence of government regulation on camera production.
    \item the market system answering the question of "where" cameras should be produced.
\end{enumerate}
\item If your nominal income rises 4 percent and your real income falls 1 percent by how much did the price level change?
\begin{enumerate}[label=(\alph*)]
    \item 4 percent increase
    \item 3 percent increase
    \item No change in the price level
    \item 5 percent decrease
    \item 5 percent increase
\end{enumerate}
\item Which of the following statements is true concerning the population regression function (PRF) and sample regression function (SRF)?
\begin{enumerate}[label=(\alph*)]
    \item The PRF and SRF are interchangeable terms referring to the same concept.
    \item The PRF is the model used for prediction, while the SRF is used for explanation
    \item The PRF is used to infer likely values of the SRF
    \item Whether the model is good can be determined by comparing the SRF and the PRF
    \item The SRF is a theoretical construct, while the PRF is an empirical construct
\end{enumerate}
\item If corn is produced in a perfectly competitive market and the government placed a price ceiling above equilibrium, which of the following would be true?
\begin{enumerate}[label=(\alph*)]
    \item Corn would become a luxury good due to the price ceiling.
    \item The government would have to subsidize corn production.
    \item The price of corn would increase dramatically.
    \item The price of corn would remain the same but the quality would decrease.
    \item The producers of corn would lose revenue due to the decreased price.
\end{enumerate}
\item What is meant by the technically efficient methodof production?
\begin{enumerate}[label=(\alph*)]
    \item A method that minimizes the time taken to produce a product
    \item A method that uses more of all inputs to increase output
    \item A method that produces the highest quantity of output
    \item A method that relies solely on the latest technology
    \item A method that achieves the fastest production speed regardless of resource usage
\end{enumerate}
\end{enumerate}

\section*{True / False (10 marks)}
\textit{Answer True or False}
\begin{enumerate}[label=\arabic*.]
\setcounter{enumi}{5}
\item True or False: Money functions as a unit of account, a store of wealth, and a standard of deferred payment in modern economies.
\item True or False: Crowding out happens when interest rates fall due to increased government borrowing, stimulating further economic growth.
\item True or False: According to Keynesian theory, changes in the money supply have significant effects on economic activity and aggregate demand.
\item True or False: An increase in interest rates in the United States typically results in the appreciation of the U.S. dollar against the euro.
\item True or False: Skylar's total benefit from eating cake is decreasing with each additional piece consumed.
\end{enumerate}

\section*{Long Form Response (20 marks)}
\textit{Answer each question with a clear and complete explanation.}
\begin{enumerate}[label=\arabic*.]
\setcounter{enumi}{10}
\item Discuss the impact of an increase in the price of pears on consumer behavior, specifically focusing on the expected changes in the quantity demanded. Analyze the underlying economic principles.
\item Discuss how it is possible for numerous firms to sell identical products within a monopolistic competition framework, and analyze the implications for pricing and consumer choice.
\end{enumerate}

\vfill
\noindent\textit{End of Paper}
\end{document}